\documentclass[../main]{subfiles}

% Chapter 3: conclusions and outlook, followed by their translation into
% Spanish (a requirement of the UAM doctoral school). The conclusions are
% organized by question (TU Delft style): a chapter lead-in, a Conclusions
% section of titled \paragraphs closed by an overview after \sectionclose,
% and an Outlook section of untitled paragraphs.

\begin{document}
\ifSubfilesClassLoaded{\setcounter{chapter}{2}}{}

\chapter{Conclusions and Outlook}
\label{ch:ch3}
\begin{otherlanguage}{french}
\begin{chquote}
Ceci est une citation d'exemple. \\
\bigskip
-- Auteur, \textit{Œuvre}
\end{chquote}
\end{otherlanguage}

% A chapter without chabstract fills the opening page with \vfill.
\vfill

\clearpage

The lead-in recalls the question that runs through the thesis.
\lipsum[1][1-6]

\section{Conclusions}
\label{sec:conclusions}

\paragraph{The first question}
\lipsum[2][1-6]

\paragraph{The second question}
\lipsum[3][1-6]

\paragraph{The third question}
\lipsum[4][1-6]

\paragraph{The fourth question}
\lipsum[5][1-6]

\sectionclose
This overview gathers the answers to the four questions.
\lipsum[6][1-3]

\section{Outlook}
\label{sec:outlook}
\lipsum[7][1-5]

\lipsum[8][1-5]

\sectionclose
\lipsum[9][1-3]


% ── Spanish translation ─────────────────────────────────────
\begin{otherlanguage}{spanish}
\chapter*{Conclusiones y perspectivas}
\addcontentsline{toc}{chapter}{\protect\numberline{}Conclusiones y perspectivas}
\label{ch:ch3_es}
% \chapter* does not set the running mark, so without this the head would
% still read "3. Conclusions and Outlook" from the chapter above.
\markboth{Conclusiones y perspectivas}{}

\vfill
\begin{disclaimer}
    Este capítulo es una traducción al español del Capítulo~\ref{ch:ch3}.\\
    \indent\foreignlanguage{english}{This chapter is a Spanish translation of Chapter~\ref{ch:ch3}.}
\end{disclaimer}

\clearpage

% Acronyms keep their English short forms after a Spanish long form written
% by hand, since \acrfull/\acrlong would print the English one. Acronym
% plurals are invariable in Spanish ("los ABS"), so \acrshort throughout.
% \section* keeps the sections out of the numbering; \markright sets the
% running head.
La introducción recuerda la pregunta que vertebra la tesis, con los estados ligados de Andreev (\acrshort{abs}, por sus siglas en inglés) como ejemplo de acrónimo.
\lipsum[1][1-6]

\section*{Conclusiones}
\markright{Conclusiones}

\paragraph{La primera pregunta}
\lipsum[2][1-6]

\paragraph{La segunda pregunta}
\lipsum[3][1-6]

\paragraph{La tercera pregunta}
\lipsum[4][1-6]

\paragraph{La cuarta pregunta}
\lipsum[5][1-6]

\sectionclose
\lipsum[6][1-3]

\section*{Perspectivas}
\markright{Perspectivas}
\lipsum[7][1-5]

\lipsum[8][1-5]

\sectionclose
\lipsum[9][1-3]
\end{otherlanguage}

\ifSubfilesClassLoaded{%
\newrefcontext[labelprefix=]
\printbibliography
}{}
\end{document}
